\subsection{Limits and colimits}
\begin{definition}
    Let $\catj$ be a small category. A \textbf{diagram $D$ of shape $\catj$} is
    any functor $D: \catj \to \catc$. Define $\Delta:\catc\to\funcat{\catj}{\catc}$
    which sends $A$ to the constant functor $\Delta A$ that 
    assigns all $j\in\ob{\catj}$ to $A$ and every morphism in $\catj$ 
    to the identity morphism $1_A$. Given $A\in\ob{\catc}$ a \textbf{cone with vertex $A$}
    is a pair $(A,\lambda)$, where $\lambda:\Delta A\to D$ is a natural transformation.
    A \textbf{universal cone} is a cone $(A,\lambda)$ so that for any other
    cone $(B,\mu)$ there exists a unique morphism $f: B\to A$ such that
    $\forall j\in\catj : \lambda_j \circ f = \mu_j$.
\end{definition}
\begin{lemma}
    Universal cone of a diagram $D$ is unique up to a unique isomorphism, if it exists.
\end{lemma}
\begin{proof}[Proof]
    
\end{proof}
\begin{definition}
    The \textbf{limit of a diagram $D$} is its universal cone, if it exists.
\end{definition}
\begin{lemma}[Pullback lemma]
    Consider this commutative diagram
    \[\begin{tikzcd}
	\bullet && \bullet && \bullet \\
	\\
	\bullet && \bullet && \bullet
	\arrow["{f_1}"', from=1-1, to=1-3]
	\arrow["{f_2}", from=1-1, to=3-1]
	\arrow["{f_3}"', from=1-3, to=1-5]
	\arrow["{p_1}", from=1-3, to=3-3]
	\arrow["{q_1}"', from=1-5, to=3-5]
	\arrow["{p_2}", from=3-1, to=3-3]
	\arrow["{q_2}", from=3-3, to=3-5]
\end{tikzcd}\]
    assuming that the right square is a pullback, then 
    the left square is a pullback iff the outer rectangle
    is a pullback.
\end{lemma}
\begin{proof}[Proof]
    Assume that the outer rectangle is a pullback and 
    consider any cone $(U,s_1,s_2)$ like so 
    \[\begin{tikzcd}
	U \\
	& \bullet && \bullet && \bullet \\
	\\
	& \bullet && \bullet && \bullet
	\arrow["{s_1}"', curve={height=-6pt}, from=1-1, to=2-4]
	\arrow["{s_2}"', curve={height=6pt}, from=1-1, to=4-2]
	\arrow["{f_1}"', from=2-2, to=2-4]
	\arrow["{f_2}", from=2-2, to=4-2]
	\arrow["{f_3}"', from=2-4, to=2-6]
	\arrow["{p_1}", from=2-4, to=4-4]
	\arrow["{q_1}"', from=2-6, to=4-6]
	\arrow["{p_2}", from=4-2, to=4-4]
	\arrow["{q_2}", from=4-4, to=4-6]
\end{tikzcd}\]
    since the outer rectangle is a pullback there exists a unique 
    morphism $m$ that makes the triangles in the following diagram commute
    \[\begin{tikzcd}
	U \\
	& \bullet && \bullet \\
	\\
	& \bullet && \bullet
	\arrow["m"{description}, from=1-1, to=2-2]
	\arrow["{f_3\circ s_1}", curve={height=-6pt}, from=1-1, to=2-4]
	\arrow["{s_2}"', curve={height=6pt}, from=1-1, to=4-2]
	\arrow["{f_3\circ f_1}"', from=2-2, to=2-4]
	\arrow["{f_2}", from=2-2, to=4-2]
	\arrow["{q_1}"', from=2-4, to=4-4]
	\arrow["{q_2\circ p_2}"', from=4-2, to=4-4]
\end{tikzcd}\] 
    because $(U,f_3\circ s_1, s_2)$ is a cone above the outer rectangle. Now we have
    to prove that $f_1\circ m = s_1$
    \[\begin{tikzcd}
	U &&&&&&& U \\
	& \bullet && \bullet && \bullet &&& \bullet & \bullet \\
	&&&&&&&& \bullet & \bullet \\
	& \bullet && \bullet && \bullet
	\arrow["m"{description}, from=1-1, to=2-2]
	\arrow["{{s_1}}"', curve={height=-6pt}, from=1-1, to=2-4]
	\arrow["{{s_2}}"', curve={height=6pt}, from=1-1, to=4-2]
	\arrow["{f_1\circ m}"{description}, curve={height=6pt}, from=1-8, to=2-9]
	\arrow["{s_1}"{description}, curve={height=-6pt}, from=1-8, to=2-9]
	\arrow["{f_3\circ s_1}"{description}, curve={height=-18pt}, from=1-8, to=2-10]
	\arrow["{p_1\circ f_1\circ m}"{description}, curve={height=18pt}, from=1-8, to=3-9]
	\arrow["{{f_1}}"', from=2-2, to=2-4]
	\arrow["{{f_2}}", from=2-2, to=4-2]
	\arrow["{{f_3}}"', from=2-4, to=2-6]
	\arrow["{{p_1}}", from=2-4, to=4-4]
	\arrow["{{q_1}}"', from=2-6, to=4-6]
	\arrow["{f_3}", from=2-9, to=2-10]
	\arrow["{p_1}", from=2-9, to=3-9]
	\arrow["{q_1}", from=2-10, to=3-10]
	\arrow["{q_2}"', from=3-9, to=3-10]
	\arrow["{{p_2}}", from=4-2, to=4-4]
	\arrow["{{q_2}}", from=4-4, to=4-6]
\end{tikzcd}\]
    By a diagram chase we verify, that the outer square on the right side 
    really commutes. From the universal property of pullbacks we get that 
    $f_1\circ m = s_1$ as required. Now to show that
    $m$ is unique consider the commutative diagram
    \[\begin{tikzcd}
	U \\
	& \bullet && \bullet \\
	\\
	& \bullet && \bullet
	\arrow["m"{description}, shift right=2, from=1-1, to=2-2]
	\arrow["{m'}"{description}, shift left=2, from=1-1, to=2-2]
	\arrow["{s_1\circ f_3}", curve={height=-6pt}, from=1-1, to=2-4]
	\arrow["{{s_2}}"', curve={height=6pt}, from=1-1, to=4-2]
	\arrow["{f_3\circ f_1}"', from=2-2, to=2-4]
	\arrow["{{f_2}}", from=2-2, to=4-2]
	\arrow["{{q_1}}"', from=2-4, to=4-4]
	\arrow["{q_2\circ p_2}"', from=4-2, to=4-4]
\end{tikzcd}\]
    it follows by the universal property that $m=m'$, since the outer
    rectangle is a pullback. Now assume that the left square is a pullback
    and consider any cone $(U,s_1,s_2)$
\[\begin{tikzcd}
	U \\
	& \bullet && \bullet && \bullet \\
	\\
	& \bullet && \bullet && \bullet
	\arrow["{s_1}"{description}, curve={height=-6pt}, from=1-1, to=2-6]
	\arrow["{{s_2}}"', curve={height=6pt}, from=1-1, to=4-2]
	\arrow["{{f_1}}"', from=2-2, to=2-4]
	\arrow["{{f_2}}", from=2-2, to=4-2]
	\arrow["{{f_3}}"', from=2-4, to=2-6]
	\arrow["{{p_1}}", from=2-4, to=4-4]
	\arrow["{{q_1}}"', from=2-6, to=4-6]
	\arrow["{{p_2}}", from=4-2, to=4-4]
	\arrow["{{q_2}}", from=4-4, to=4-6]
\end{tikzcd}\]
    Since $(U,s_1, p_2\circ s_2)$ is a cone and the right square is a pullback,
    there exists unique $n$ such that $f_3\circ n = s_1$ and $p_1\circ n = p_2\circ s_2$,
    but then $(U,s_2,n)$ is a cone on the left pullback square, so there exists a unique
    $m$ such that $f_1\circ m = n$ and $f_2\circ m = s_2$ like so 
\[\begin{tikzcd}
	U \\
	& \bullet && \bullet && \bullet \\
	\\
	& \bullet && \bullet && \bullet
	\arrow["m"{description}, from=1-1, to=2-2]
	\arrow["n"{description}, curve={height=-6pt}, from=1-1, to=2-4]
	\arrow["{s_1}"{description}, curve={height=-12pt}, from=1-1, to=2-6]
	\arrow["{{s_2}}"', curve={height=6pt}, from=1-1, to=4-2]
	\arrow["{{f_1}}"', from=2-2, to=2-4]
	\arrow["{{f_2}}", from=2-2, to=4-2]
	\arrow["{{f_3}}"', from=2-4, to=2-6]
	\arrow["{{p_1}}", from=2-4, to=4-4]
	\arrow["{{q_1}}"', from=2-6, to=4-6]
	\arrow["{{p_2}}", from=4-2, to=4-4]
	\arrow["{{q_2}}", from=4-4, to=4-6]
\end{tikzcd}\]
    simple diagram chase shows that $m$ completes the cone $(U,s_1,s_2)$.
    It remains to show that $m$ is unique. Let $m'$ be any other completion
    of the cone $(U,s_1,s_2)$
    \[\begin{tikzcd}
	U \\
	& \bullet && \bullet && \bullet \\
	\\
	& \bullet && \bullet && \bullet
	\arrow["{m'}"{description}, from=1-1, to=2-2]
	\arrow["n"{description}, curve={height=-6pt}, from=1-1, to=2-4]
	\arrow["{s_1}"{description}, curve={height=-12pt}, from=1-1, to=2-6]
	\arrow["{{s_2}}"', curve={height=6pt}, from=1-1, to=4-2]
	\arrow["{{f_1}}"', from=2-2, to=2-4]
	\arrow["{{f_2}}", from=2-2, to=4-2]
	\arrow["{{f_3}}"', from=2-4, to=2-6]
	\arrow["{{p_1}}", from=2-4, to=4-4]
	\arrow["{{q_1}}"', from=2-6, to=4-6]
	\arrow["{{p_2}}", from=4-2, to=4-4]
	\arrow["{{q_2}}", from=4-4, to=4-6]
\end{tikzcd}\]
    if we show $f_1\circ m' = n$ then by the universal property we would have
    that $m=m'$. By a diagram chase we see that $(U, s_1, p_2\circ s_2)$ is 
    a cone completed by $n$ and $f_1\circ m'$ of the right pullback square
    \[\begin{tikzcd}
	U \\
	&&& \bullet && \bullet \\
	\\
	&&& \bullet && \bullet
	\arrow["n"{description}, curve={height=-6pt}, from=1-1, to=2-4]
	\arrow["{f_1\circ m'}"{description}, curve={height=6pt}, from=1-1, to=2-4]
	\arrow["{s_1}"{description}, curve={height=-12pt}, from=1-1, to=2-6]
	\arrow["{p_2\circ s_2}"{description}, curve={height=12pt}, from=1-1, to=4-4]
	\arrow["{{f_3}}"', from=2-4, to=2-6]
	\arrow["{{p_1}}", from=2-4, to=4-4]
	\arrow["{{q_1}}"', from=2-6, to=4-6]
	\arrow["{{q_2}}", from=4-4, to=4-6]
\end{tikzcd}\]
therefore from universality $f_1\circ m' = n$ and again by universality of the left pullback $m=m'$.
\end{proof}
\begin{proposition}
    Given a diagram 
    \[\begin{tikzcd}
	E & A & B
	\arrow["e", from=1-1, to=1-2]
	\arrow["{f_1}", shift left, from=1-2, to=1-3]
	\arrow["{f_2}"', shift right, from=1-2, to=1-3]
\end{tikzcd}\]
    such that $(E,e)$ is the equaliser of $f_1$ and $f_2$, then
    $e$ is a monomorphism.
\end{proposition}
\begin{proof}[Proof]
    Consider a pair of morphisms $g_1,g_2: A'\to E$ such that $eg_1=eg_2$, then
    \[\begin{tikzcd}
	E & A & B \\
	{A'}
	\arrow["e", from=1-1, to=1-2]
	\arrow["{f_1}", shift left, from=1-2, to=1-3]
	\arrow["{f_2}"', shift right, from=1-2, to=1-3]
	\arrow["{\exists !}", dashed, from=2-1, to=1-1]
	\arrow["{eg_1=eg_2}"', from=2-1, to=1-2]
\end{tikzcd}\]
    therefore $g_1=g_2$ since they both also make the triangle on the left commute.
\end{proof}
\begin{proposition}
    Consider the following pullback diagram
    \[\begin{tikzcd}
	P && A \\
	\\
	B && C
	\arrow["{m_1}", from=1-1, to=1-3]
	\arrow["{f_1}"', from=1-1, to=3-1]
	\arrow["\lrcorner"{anchor=center, pos=0.125}, draw=none, from=1-1, to=3-3]
	\arrow["{f_2}", from=1-3, to=3-3]
	\arrow["{m_2}"', hook, from=3-1, to=3-3]
\end{tikzcd}\]
    then $m_1$ is a monomorphism.
\end{proposition}
\begin{proof}[Proof]
    Suppose that $g_1$ and $g_2$ are such that $m_1 g_1 = m_1 g_2$
    then by a diagram chase we can easily show that $f_1 g_2 = f_1 g_1$ 
    so that we have
    \[\begin{tikzcd}
	X \\
	& P && A \\
	\\
	& B && C
	\arrow["{g_1}"{description}, curve={height=-6pt}, from=1-1, to=2-2]
	\arrow["{g_2}"{description}, curve={height=6pt}, from=1-1, to=2-2]
	\arrow["{m_1\circ g_1=m_1\circ g_2}", curve={height=-6pt}, from=1-1, to=2-4]
	\arrow["{f_1\circ g_2 = f_1\circ g_1}"', curve={height=12pt}, from=1-1, to=4-2]
	\arrow["{m_1}", from=2-2, to=2-4]
	\arrow["{f_1}"', from=2-2, to=4-2]
	\arrow["\lrcorner"{anchor=center, pos=0.125}, draw=none, from=2-2, to=4-4]
	\arrow["{f_2}", from=2-4, to=4-4]
	\arrow["{m_2}"', hook, from=4-2, to=4-4]
\end{tikzcd}\]
    since we have a pullback square it follows that $g_1=g_2$.
\end{proof}
\begin{definition}
    Let $\catc$ be a category and choose $X\in\ob{\catc}$. We define the \textbf{slice category $\catc$ over $X$} denoted as $\catc /X$,
    whose objects are pairs $(A,f)$ where $A\in\ob{\catc}$, $f: A \to X$ and whose morphisms 
    $g: (A,f_1) \to (B,f_2)$ are just morphisms $g:A\to B$ in $\catc$ such that the following diagram commutes
\[\begin{tikzcd}
	A && B \\
	& X
	\arrow["g", from=1-1, to=1-3]
	\arrow["{f_1}", from=1-1, to=2-2]
	\arrow["{f_2}"', from=1-3, to=2-2]
\end{tikzcd}\]
    Notice that composition of two morphisms $\catc /X$ yields again a morphisms in $\catc /X$
    since when the two triangles in the left diagram below commute then the entire square commutes as well
    and that by the very definition of identity morphism the triangle on the right side below commutes 
\[\begin{tikzcd}
	& B &&& A && A \\
	A && C \\
	& X &&&& X
	\arrow["{g_2}", from=1-2, to=2-3]
	\arrow["{f_2}"{description}, from=1-2, to=3-2]
	\arrow["{1_A}", from=1-5, to=1-7]
	\arrow["{f_1}"', from=1-5, to=3-6]
	\arrow["{f_1}", from=1-7, to=3-6]
	\arrow["{g_1}", from=2-1, to=1-2]
	\arrow["{f_1}"', from=2-1, to=3-2]
	\arrow["{f_3}", from=2-3, to=3-2]
\end{tikzcd}\]
\end{definition}
\begin{definition}
    Given two objects $A,B$ their coproduct is an object $U$ with 
    two morphisms $\iota_A: A \to U$ and $\iota_B: B \to U$, such
    that for any other triple $(U^\prime,f: A\to U^\prime,g: A\to U^\prime)$
    there exists a unique morphism $f+g: U\to U^\prime$, such that the following diagram commutes
\[\begin{tikzcd}
	& {U^\prime} \\
	\\
	A & U & B
	\arrow["f", from=3-1, to=1-2]
	\arrow["{\iota_A}", from=3-1, to=3-2]
	\arrow["{f+g}"{description}, from=3-2, to=1-2]
	\arrow["g"', from=3-3, to=1-2]
	\arrow["{\iota_B}"', from=3-3, to=3-2]
\end{tikzcd}\]
\end{definition}
\begin{proposition}
    Coproduct $U$ is unique up to an isomorphism.
\end{proposition}